\chapter{Coverage, Descent, and Multiple Views}
\label{ch:descent}

Chapters~\ref{ch:sheaves} and~\ref{ch:cohomology} took time as the base space. Production also works with a cover of the \emph{scene} by many cameras. This chapter treats the second kind of cover, relates it to the problem of contextuality studied in quantum foundations, and states a result on when a family of local reconstructions determines a global one.

\section{Multi-camera coverage}

A scene is staged in a region $X\subseteq\R^3$ over a time interval. Cameras $i\in I$ each see an open subregion $U_i\subseteq X$, their \emph{field of view}, and each camera's footage supports a local reconstruction of the scene state on $U_i$. Let $F$ be the sheaf of scene states on $X$, with $F(U)$ the set of consistent descriptions of the state of the region $U$, and suppose each camera produces a section $s_i\in F(U_i)$.

\begin{definition}[Coverage]
The family of cameras provides \emph{coverage} if $\bigcup_iU_i=X$. It provides \emph{consistent coverage} if it provides coverage and $\restr{s_i}{U_{ij}}=\restr{s_j}{U_{ij}}$ for all $i,j$.
\end{definition}

By the sheaf axioms, consistent coverage determines a unique global state. Production practice contains a method for checking the hypotheses: continuity supervisors compare overlapping views for matching props, costume position, and the like, which is a check of the restriction equalities on $U_{ij}$.

\begin{proposition}\label{prop:cover}
If coverage is consistent, the global scene state $s\in F(X)$ is uniquely determined by $(s_i)$. If coverage is not consistent, there exist indices $i,j$ with $\restr{s_i}{U_{ij}}\ne\restr{s_j}{U_{ij}}$, and no global state restricts to both $s_i$ and $s_j$.
\end{proposition}

\begin{proof}
The first statement is the gluing and locality axioms. For the second, suppose $\restr{s_i}{U_{ij}}\ne\restr{s_j}{U_{ij}}$ in $F(U_{ij})$ for some $i,j$. A global section $s\in F(X)$ restricting to both $s_i$ and $s_j$ would satisfy $\restr{s_i}{U_{ij}}=\restr{(\restr{s}{U_i})}{U_{ij}}=\restr{s}{U_{ij}}=\restr{s_j}{U_{ij}}$ by functoriality of restriction, a contradiction.
\end{proof}

The proposition has a purely practical counterpart: an inconsistency visible in the overlap of two cameras is a certificate that the footage cannot all depict the same event.

\section{Contextuality and empirical models}

The pattern of local consistency without global consistency is not specific to film. It is the structure of contextuality in quantum mechanics, and the sheaf-theoretic account of Abramsky and Brandenburger~\cite{abramsky2011} gives a precise analogue. In that setting, one has a set of measurements, a cover by the maximal sets of jointly performable measurements, and local probability distributions on each such context that agree on overlaps. The empirical model is \emph{non-contextual} if it extends to a global distribution and \emph{contextual} otherwise. Local agreement is a section of a compatible family, and the failure to extend is the failure of gluing.

(The analogy uses a presheaf of local distributions or accounts, which need not be a sheaf; this differs from the sheaf $F$ of scene states above, for which gluing holds by assumption.) The filmic analogue is a collection of viewpoints, each internally coherent and each agreeing with the others where they overlap, with no single world state accounting for all of them. This is a technical rendering of the claim that some narratives are told from perspectives that cannot be reconciled in one objective account. Whether a specific film, such as \emph{Last Year at Marienbad} (1961), has this structure is a matter of reading, and the model allows one to state the claim exactly: that the family of recalled scenes is locally compatible and admits no global section.

\begin{remark}
The analogy with quantum contextuality should be taken only as far as the mathematics goes. The cover and sheaf condition are the common structure. Nothing here implies that films exhibit quantum behavior, and nothing in the physics of contextuality is claimed for narrative.
\end{remark}

\section{Descent}

A more refined statement relates the sheaf condition to a universal property. Given a cover $\mathcal{U}$ of $X$, the sheaf condition says that $F(X)$ is the equalizer of two maps
\[
\prod_iF(U_i)\rightrightarrows\prod_{i,j}F(U_{ij}),
\]
which send a family $(s_i)$ to $(\restr{s_i}{U_{ij}})$ and $(\restr{s_j}{U_{ij}})$ respectively. In the language of descent, a global section is the same as a family of local sections together with the equality of their two restrictions on all overlaps~\cite{maclane1992}.

When the local data are not sets but categories, as when each camera's footage supports a category of possible reconstructions rather than a single one, equality on overlaps is replaced by a specified isomorphism, and these isomorphisms must satisfy a cocycle condition on triple overlaps. The resulting structure is a \emph{descent datum}, and its effectiveness (the existence of a global object inducing it) is a stronger condition than gluing of sets. The cocycle condition is the non-abelian form of the \v{C}ech 1-cocycle condition of the previous chapter. Obstructions to the existence of descent data at all lie one level higher, in second cohomology (gerbes), which we do not pursue. (Strictly, the local data then form a stack, with restriction defined up to canonical isomorphism.)

\begin{definition}[Descent datum for coverage]
A descent datum for cameras $(U_i)$ with local reconstruction categories $F(U_i)$ consists of objects $s_i\in F(U_i)$ and isomorphisms $\phi_{ij}:\restr{s_i}{U_{ij}}\to\restr{s_j}{U_{ij}}$ such that $\restr{\phi_{jk}}{U_{ijk}}\circ\restr{\phi_{ij}}{U_{ijk}}=\restr{\phi_{ik}}{U_{ijk}}$ on triple overlaps. The datum is \emph{effective} if there exists a global $s\in F(X)$ with isomorphisms $\restr{s}{U_i}\cong s_i$ compatible with the $\phi_{ij}$.
\end{definition}

An editor assembling a scene from several cameras supplies the isomorphisms $\phi_{ij}$ implicitly by identifying props and actors across views. The cocycle condition is the requirement that this identification be transitive: if prop $a$ in camera $1$ is prop $b$ in camera $2$ and prop $b$ is prop $c$ in camera $3$, then $a$ is $c$. A violation means the identifications do not form a descent datum, a continuity error of the \emph{identity} kind. Effectiveness is a further property of a datum that does satisfy the cocycle condition, namely that some single global scene accounts for it.

\section{Bridge: when is gluing the right demand?}

Suppose the forged-letter scene is covered by three cameras. Each camera yields local sections of several separately recorded sheaves: the position of objects, the clothing of the characters, the lighting, what each character knows, and the timing of events. Comparison on overlaps proceeds sheaf by sheaf. Camera one shows the letter on the left of the table at $t=30$, and camera two, whose view overlaps camera one's, shows it on the right at $t=38$. Camera three shows $B$'s sleeve rolled up at $t=44$ and rolled down at $t=46$. These readings are at different instants, so they lie on no common overlap in time; they become a disagreement only under the assumption that the letter's position and the sleeve are constant over the interval, and that assumption is what a gluing demand makes.

A mismatch of this kind has several possible explanations, and the mathematics does not choose among them. It may be a production error, of the sort a continuity supervisor exists to catch. It may be a deliberate ambiguity, where the film wishes the audience to be unsure of the scene's physical state. It may be a subjective section, such as a recollection, for which the demand for agreement with another view was never appropriate. Or it may show that the model is inadequate: the letter may have been moved during the pause between $t=31$ and $t=38$, in which case the two position readings, taken before and after the pause, do not conflict once the history contains the event of its being moved, and the apparent contradiction came from asking for a single state where the story has two.

This suggests when gluing is an appropriate demand. The requirement that local sections agree on overlaps is justified for a sheaf that models objective world state, over a stretch of time that the film presents as continuous and objective. It is not justified across a cut to another time, or for a section marked as someone's recollection, which belongs in a separate sheaf indexed by the person who recollects. Only after the sheaf and the stretch have been chosen is a failure to glue evidence of anything. The last explanation above also shows how the apparatus of Chapter~\ref{ch:layers} enters: a section that fails to glue as a state may glue once the history is enlarged by the missing event.

\section*{Exercises}

\begin{exercise}
Let $X=[0,3]$ with $U_1=[0,2)$, $U_2=(1,3]$. Give an example of a sheaf $F$ and local sections $s_1,s_2$ that agree on $U_{12}$ but whose glued section fails to belong to a given subpresheaf $F'\subset F$ of ``allowed'' states. What does this say about checking constraints locally?
\end{exercise}

\begin{exercise}
Verify that the cocycle condition of a descent datum is automatic when the isomorphisms $\phi_{ij}$ arise from a global object, and that it is the only extra condition needed at the triple-overlap level. (Show also that $\phi_{ii}=\mathrm{id}$ follows from the cocycle condition.)
\end{exercise}

\begin{exercise}
Describe a three-camera scene in which each pair of cameras is consistent but the three together are not, in the sense that the transitivity of prop identification fails. Is such a configuration ruled out by the sheaf axiom for sets?
\end{exercise}
